%% The following is a directive for TeXShop to indicate the main file
%%!TEX root = diss.tex

%=======================================
\chapter{The Obstinate Kernel Method}
\label{ch:OKM}
%=======================================

We have another tool at our disposal, the Obstinate Kernel method, which was used in the interacting case in \cite{tabbara2014exact}. The idea is similar to the weighted algebraic kernel method in that we produce more functional equations from the master equation via group actions $g \in G$. The difference is that here we also set one of $x_1 = \sigma_1$, $x_2 = \sigma_2$ equal to a root of $K$ in order to eliminate the $KQ$ term from the equations. This leaves equations that only include the other variable. It turns out that some combinations of substitutions and group operations do not lead to valid functional equations. This happens when one of the $Q_k^0$ terms is not a formal power series in $t$. In all the cases with finite groups of order 4, 6 and 8 we find that there are 2 roots of $K$ for each $x_k$ and that one root is in $O(t)$ and the other in $O(1/t)$. Generally, it works out better if we use the $O(t)$ roots because that at least guarantees that the $Q_k^0$ terms remain as formal power series in $t$. However, sometimes (page 11 of \cite{tabbara2014exact}) even the $O(1/t)$ term may lead to a valid power series. 

It may not be necessary, but it will be sufficient if a combination of a substitution $x_k = \sigma_k$ and a group action $g_L$ together give a power series in $t$ when applied to single variables like $x_1$ or $x_2$. For example, we would hope that $g_L(x_j)|_{x_k=\sigma_k}$ is a power series in $t$ for some $i,j,L$. Again, this is not a necessary condition, but it will be sufficient to show that a lot of the cases involved produce enough different equations to allow us to go through the Obstinate Kernel method. For those cases where we find less equations, it is likely that we will be able to show that the $Q_k^0$ terms are power series in $t$ via another method like in (two friendly walks).

\begin{definition} A function $f(\row xd;t)$ is admissible if it is a power series in $t$. If $\sigma_k$ is a root of $K$, we say $f$ is $k$-admissible if $f|_{x_k=\sigma_k}$ is admissible. A set of functions is admissible if each element is admissible.
\end{definition}

\begin{prop} Let $\mc{M}$ be a lattice walk model with small nonrepeating (?) steps and group $G$ of order 4. The actions $g \in G$ that preserve the admissibility are given in the following table.
\end{prop}

When $|G| = 6$ we have a richer situation.

\begin{prop} Let $\mc{M}$ be a lattice walk model with small nonrepeating steps and group $G$ of order 4, 6 or 8. Models with the same group have the same admissibility properties with respect to admissible when it comes to $x_1$, $x_2$ or the pair $(x_1,x_2)$.
\end{prop}
\begin{proof} We do not know of a conceptual proof of this fact, so we performed the calculations on a computer. They are trivial computationally speaking. For a given model $\mc{M}$, we construct its $S$-polynomial, which can be used to find the admissible roots $\sigma_k$ for $k =1,2$. We can use $S$ to express the generators $g_k$ and their action on the pair $(x_1,x_2)$. Once we have the generators we form the group $G$. Next, to check the admissibility of $f\in \{x_1$, $x_2$, $(x_1,x_2)\}$ all we need to do is find $|\orb_G(f)|$ and see if substituting $x_k = \sigma_k$ in the resulting expressions leads to admissible functions. The calculations can be done by hand, but they are tedious, so we have employed a computer to do the legwork. It becomes a trivial matter to check the admissibility properties for the various models and compare the results, summarized in the following tables.
\end{proof}

\begin{table}[H]
\caption{$k$-admissibility of functions when $|G| = 4$.}
\label{tbl:Admissibility4}
\centering
\begin{tabular}{cccc}
\toprule
	Type & $x_1$  & $x_2$ & $(x_1,x_2)$ \\
\midrule
	$1$-admissibility & $1$, $ g_2$  & all $g\in G$ & $1$, $ g_2$\\
	$2$-admissibility & $1$, $ g_1$  & all $g\in G$ & $1$, $ g_1$\\
\bottomrule
\end{tabular}
\end{table}

\begin{table}[H]
\caption{$k$-admissibility of functions for $S11$, $S12$, $S13$. Here $|G| = 6$.}
\label{tbl:Admissibility61}
\centering
\begin{tabular}{cccc}
\toprule
	Type & $x_1$  & $x_2$ & $(x_1,x_2)$ \\
\midrule
	$1$-admissibility & $1$, $ g_{2} $, $ g_{21} $, $ g_{121}$ & $1$, $ g_{1} $, $ g_{2} $, $ g_{21}$ & $1$, $ g_{2} $, $ g_{21} $ \\
	$2$-admissibility & $1$, $ g_{1} $, $ g_{2} $, $ g_{12}$ & $1$, $ g_{1} $, $ g_{12} $, $ g_{121}$ & $1$, $ g_{1} $, $ g_{12}$ \\
\bottomrule
\end{tabular}
\end{table}

\begin{table}[H]
\caption{$k$-admissibility of functions for $S21$, $S22$, $S33$. Here $|G| = 6$.}
\label{tbl:Admissibility62}
\centering
\begin{tabular}{cccc}
\toprule
	Type & $x_1$  & $x_2$ & $(x_1,x_2)$ \\
\midrule
	$1$-admissibility & $1$, $ g_{2} $, $ g_{21} $, $ g_{121}$ & $1$, $ g_{1} $, $ g_{12} $, $ g_{121}$ & $1$, $ g_{121}$ \\
	$2$-admissibility & $1$, $ g_{2} $, $ g_{21} $, $ g_{121}$ & $1$, $ g_{1} $, $ g_{12} $, $ g_{121}$ & $1$, $ g_{121}$ \\
\bottomrule
\end{tabular}
\end{table}

\begin{table}[H]
\caption{$k$-admissibility of functions for $S1$. Here $|G| = 8$.}
\label{tbl:Admissibility81}
\centering
\begin{tabular}{cccc}
\toprule
	Type & $x_1$  & $x_2$ & $(x_1,x_2)$ \\
\midrule
	$1$-admissibility & $1$, $ g_{2} $, $ g_{21} $, $ g_{212}$ & all $g\in G$ & $1$, $ g_{2} $, $ g_{21} $, $ g_{212}$ \\
	$2$-admissibility & $1$, $ g_{1} $, $ g_{2} $, $ g_{12}$, $ g_{121}$, $ g_{2121}$ & $1$, $ g_{1} $, $ g_{12} $, $ g_{121}$ & $1$, $ g_{1} $, $ g_{12} $, $ g_{121}$ \\
\bottomrule
\end{tabular}
\end{table}

\begin{table}[H]
\caption{$k$-admissibility of functions for $S2$. Here $|G| = 8$.}
\label{tbl:Admissibility82}
\centering
\begin{tabular}{cccc}
\toprule
	Type & $x_1$  & $x_2$ & $(x_1,x_2)$ \\
\midrule
	$1$-admissibility & $1$, $ g_{2} $, $ g_{21} $, $ g_{212}$ & all $g\in G$ & $1$, $ g_{2} $, $ g_{21} $, $ g_{212}$ \\
	$2$-admissibility & $1$, $ g_{2} $, $ g_{21} $, $ g_{121}$, $ g_{212}$, $ g_{2121}$ & $1$, $ g_{1} $, $ g_{12} $, $ g_{121}$ & $1$, $ g_{121}$ \\
\bottomrule
\end{tabular}
\end{table}

What this classification tells us is that we can easily treat all $|G| = 4$ models with the same method. When $|G| = 6$ we either have to do some more work in the $S2k$ cases to show that the functions are admissible with respect to one more group operation, or just deal with them as they are. The situation is similar in the $|G| = 8$ case. Now suppose that we have a set of valid equations for a substitution $x_k = \sigma_k$. They are a linear system in the variable quantities $\Qf{0,0}$, $ Q_1^0$, $ Q_2^0$ and their orbits under the set of valid group actions. We hope for a high level of redundancy in these equations. Then we can eliminate some variables by taking a weighted sum over the valid equations. That is, we attain a weighted Orbit Sum operator that only has a few group operations, perhaps. We also have to keep in mind that we are also making the substitution $x_k = \sigma_k$.

We will perform the first steps in the analysis of the Obstinate Kernel method for each situation that has $|G|/2$ $k$-admissible elements since that seems to work.
%=======================================
\section{Group of Order Four}
%=======================================
The analysis in this case is the same whether we work with the $\sigma_1$ or $\sigma_2$ substitutions. We will just work with the $\sigma_1$-substitution and group actions $1$, $ g_2$, where we let $g_2(F) = F^\ast$ to keep the notation cleaner. The functional equations corresponding the the group actions are
\begin{align}
	KQ - K_1 Q_1^0 - K_2 Q_2^0 + K_{12}\Qf{0,0} &= \alpha\,,\\
	KQ^\ast \exxx - K_1^\ast Q_1^{0,\ast}\exxxx- K_2^\ast \!Q_2^0 - K_{12}^\ast \Qf{0,0} &= \alpha\,.
\end{align}
We will eliminate $Q_2^0$ because once we make the $\sigma_1$-substitution, we would have a tough time trying to expand $Q_2^0$ into series. We do this by taking $K_2^\ast$ times the first equation and subtract $K_2$ times the second equation. This leads to the functional equation
\begin{align}
	K_2^\ast KQ - K_2 K Q^\ast- (K_2^\ast K_1 Q_1^0 - K_2 K_1^\ast Q_1^{0,\ast}) + (K_2^\ast K_{12}-K_2 K_{12}^\ast)\Qf{0,0} - \alpha(K_2^\ast - K_2) = 0
\end{align}
where we can set $x_1 = \sigma_1$ to eliminate the $K$ term. We overload notation by not changing any of the symbols. From now on we will know that $x_1 = \sigma_1$ in the equations that follow. We have
\begin{align}
	(K_2^\ast K_1 Q_1^0 - K_2 K_1^\ast Q_1^{0,\ast}) - (K_2^\ast K_{12}-K_2 K_{12}^\ast)\Qf{0,0} + \alpha(K_2^\ast - K_2) = 0.
\end{align}
We can also think of this process as applying a weighted orbit sum over group actions that preserve the admissibility of functions in the master equation. The operator is
\begin{align}
	D_2 = g_2(K_2) - K_2 g_2 = (1-g_2)K_2^\ast
\end{align}
where the weights are given by pre-multiplying an expression by $K_2^\ast$ before applying the operator $\Delta_2 = 1-g_2$. The equation can be written as
\begin{align}\begin{split}
	0 &= D_2(KQ - K_1 Q_1^0 - K_2 Q_2^0 + K_{12}\Qf{0,0})\\
	0 &= D_2(K_1Q_1^0) - D_2(K_{12})\Qf{0,0} + D(\alpha)	    
\end{split}\end{align}
where the second one follow from substituting $x_1 = \sigma_1$.
The situation is similar to the Algebraic Kernel method except that the operators there retain the $KQ$ term and its orbit, while they eliminate the $Q_1^0$ and $Q_2^0$ terms and their orbits.
\begin{definition} (?) For a concatenated list $L = L_1L_2\ldots$ of positive integer indices, let $g_L = g_{L_1}g_{L_2}\cdots$ whenever this group action makes sense. That is, either the list is finite or there is some form of convergence when $L$ is infinite.
\end{definition}
\begin{definition} Let $L$ be a concatenated list. Definte $\Delta_L = 1 - g_L$. Now let $X$ be some function. The operator $D_L[X]$ is defined by its action on functions $F$ as
\begin{align}
	D_L[X](F) =\Delta_L(XF) = (1-g_L)(XF) = XF-g_L(XF)\,.
\end{align}
We will drop the $X$ from the notation if it is clear from context. In particular, $D_L[1] = \Delta_L$.
\end{definition}
For example, $D_2 = D_2[g_2(K_2)]$.
\begin{example} If we use $g_L(X)$ instead of $X$ in the above definition, we find that
\begin{align}
	D_L[g_L(X)] = g_L(X) - Xg_L
\end{align}
and this is essentially an Ore operator. In our analysis we will see these somewhat frequently.
\end{example}
%=======================================
\subsection{Some Calculations}
%=======================================
Before we go onto groups of order 6 we can expand some of the terms in the functional equation. First
\begin{align}
\begin{split}
	D_2(\alpha) &= \alpha(K_2^\ast -K_2) = \\
			&= \alpha((c_2- tx_2\ol \mu_2 S_2^{\um 1})-(c_2- t\ol x_2 S_2^{\um 1}))\\
			&= \alpha t(\ol x_2 - x_2\ol \mu_2(\sigma_1))S_2^{\um 1}(\sigma_1)\,.
\end{split}
\end{align}
Then
\begin{align}
\begin{split}
	D_2(K_{12}) &=K_2^\ast K_{12}- K_2 K_{12}^\ast = \\
			&=(c_2- tx_2\ol \mu_2 S_2^{\um 1})(c_1c_2-t\ol \sigma_1\ol x_2\Sf{\um 1,\um 1}) \\
			&\quad-(c_2- t\ol x_2 S_2^{\um 1})(c_1c_2-t\ol \sigma_1x_2\ol \mu_2\Sf{\um 1,\um 1})\\
			&=c_1c_2^2-c_1c_2tx_2\ol \mu_2S_2^{\um 1} - c_2t\ol \sigma_1\ol x_2\Sf{\um 1,\um 1}+ t^2\ol \sigma_1\ol \mu_2S_2^{\um 1}\Sf{\um 1,\um 1}\\
			&\quad - c_1c_2^2+c_1c_2t\ol x_2S_2^{\um 1} + c_2t\ol \sigma_1x_2 \ol \mu_2\Sf{\um 1,\um 1}-  t^2\ol \sigma_1\ol \mu_2S_2^{\um 1}\Sf{\um 1,\um 1}\\
			&=c_1c_2t(\ol x_2-x_2\ol \mu_2)S_2^{\um 1} - c_2t\ol \sigma_1(\ol x_2-x_2\ol \mu_2)\Sf{\um 1,\um 1}\\
			&=c_2t(\ol x_2-x_2\ol \mu_2(\sigma_1))(c_1S_2^{\um 1}(\sigma_1) - \ol \sigma_1\Sf{\um 1,\um 1})\\
\end{split}
\end{align}
If the model is symmetric in both directions, we see that $g_2 S_1^{\um 1} = S_1^{\um 1}$ and therefore, $g_2 K_1 = K_1$, and we may bring that factor outside of $D_2(K_1^\ast F) = K_1 D_2(F)$. It is worth investigating if this happens in other cases since this will make it easier to perform coefficient extraction later. We can answer this in the negative.
\begin{prop} Suppose that a model with group of order 4 has symmetry around the $x_i$ axis and $j\neq i$. Then $D_i[g_iK_j] = K_jD_i[1] = K_j \Delta_i$ in the $\sigma_j$-substitution. Also, $D_j[g_j(K_i)] = K_i D_j[1] = K_i \Delta_j$ if and only if the model is also symmetric around the $x_j$ axis. Thus, we can bring out the $K$ term out of both operators if and only if the model is highly symmetric.
\end{prop}
\begin{proof} Assume that the model is symmetric around the $x_1$ axis, that is, it is vertically symmetric. This implies that $S_1^{\um 1} = S_1^1$ and $g_1(S_2^k) = S_2^k$ for all $k$. This immediately tells us that $g_1(K_2) = K_2$ and $D_1[g_1K_2] = K_2 \Delta_1$. This remains true when we substitute $x_2 = \sigma_2$. This would also be true for the other operator if the model is symmetric around the $x_2$ axis. Thus, we only need to show the forward direction of the statement. We will show that if the model is not highly symmetric, the factoring will not happen. It is enough to show that $g_2(K_1) \neq K_1$ even after the $x_1 = \sigma_1$ substitution is made.

The solutions to $K = 0$ in the $x_1$ variable are
\begin{align}
	x_1 = \frac{(1-t S_1^0) \pm \sqrt{(1-tS_1^0)^2 - 4t^2(S_1^1)^2}}{2tS_1^1}\,.
\end{align}
If $\sigma_1 = u \pm v\sqrt{w}$, then it turns out that $\ol \sigma_1 = u \mp v\sqrt{w}$ due to the symmetry and
\begin{align}
	s = \sigma_1 + \ol \sigma_1 = 2u = \frac{1-tS_1^0}{t S_1^1}\,.
\end{align}
Now
\begin{align}
	g_2 K_1 = c_2 - t \ol x_1 g_2(S_1^{\um 1}(x_2)) = c_2 - t\ol x_1 S_1^{\um 1}(\ol x_2\mu_2(x_1))
\end{align}
and $g_2 K_1 = K_1$ if and only if
\begin{align}
	S_1^{\um 1}(\ol x_2\mu_2(x_1)) = S_1^{\um 1}(x_2)
\end{align}
upon making the substitution $x_1 = \sigma_1$. We will first compute $\mu_2(\sigma_1)$. To that end we have
\begin{align}
	S_2^k(\sigma_1) = \Sf{0,k}+ \Sf{1,k}(\ol \sigma_1+ \sigma_1) = \Sf{0,k}+ \Sf{1,k}\frac{1-tS_1^0}{t S_1^1}
\end{align}
and
\begin{align}
	\mu_2(\sigma_1) = \frac{S_2^{\um 1}}{S_2^1} = \frac{ \Sf{0,\um 1}+ \Sf{1,\um 1}(1-tS_1^0)/(t S_1^1)}{\Sf{0,1}+ \Sf{1,1}(1-tS_1^0)/(t S_1^1)} = \frac{\Sf{1, \um 1}+t\left(\Sf{0,\um 1}S_1^1-\Sf{1,\um 1}S_1^0\right)}{\Sf{1, 1}+t\left(\Sf{0, 1}S_1^1-\Sf{1, 1}S_1^0\right)}\,.
\end{align}
This is a rational function in $x_2$. Using software, we can easily compute this expression and further compute $S_1^{\um 1}(\ol x_2\mu_2(x_1))$ to show that we only have equality when the model is highly-symmetric.

\end{proof}
In a later section (ref) we will investigate the highly-symmetric models in more detail.
%=======================================
\section{Group of Order Six}
%=======================================
We can perform the Obstinate Kernel method on the three models in the first group. They are as follows.
\begin{figure}[H]
    \centering
\[
	\modelgraphB 01010001\quad \modelgraphB 10001010 \quad \modelgraphB 11011011
\]
    \caption{Models $S11$, $S12$, $S13$ with group of order 6.}
\end{figure}
The $k$-admissibility of $(x_1,x_2)$ is preserved by actions $1$, $ g_j$, $ g_{jk}$ for $j \neq k$. For consistency, we will work on the $1$-admissible case. The actions are $1$, $ g_2$, $ g_{21}$. We have the table from the Algebraic Kernel method.
\begin{table}[H]
    \caption{Reductions of group actions $g \in G$ on $Q_1^0$ and $Q_2^0$.}
    \label{tbl:tableG6}
    \centering
    \begin{tabular}{ccc}
    \toprule
    Sequence $L$ & $Q_1^0$ & $Q_2^0$ \\  
    \midrule
      $\emptyset$ & $\emptyset$ & $\emptyset$  \\
      2 &  2& $\emptyset$ \\
      21 &  2& 21  \\
    \bottomrule
    \end{tabular}
\end{table}
In this group we have $g_2(x_1,x_2) = (x_1,x_1\ol x_2)$ and $g_{21}(x_1,x_2) = (\ol x_2,x_1\ol x_2)$, so the terms we wish to eliminate are those involving an $x_1$ because the substitution $x_1 = \sigma_1$ will make it difficult to expand the $Q$ functions into series. That is, we wish to eliminate $X_2 = g_2Q_1^0$ and $Y = Q_2^0$. Let
\begin{align}
	\beta = K_{12}\Qf{0,0} - \alpha = \alpha(b_1b_2\Qf{0,0}-1)
\end{align}
so that we have equations
\begin{align}
	AX+BY = \beta\com A_2X_2+B_2Y = \beta_2\com A_{21}X+B_{21}Y_{21} = \beta_{21}\,.		
\end{align}
We want a linear combination of $\beta, \beta_2, \beta_{21}$ that will eliminate $X_2$ and $Y$. We may write
\begin{align}
	a\beta - b\beta_2+c \beta_{21} = aAX + (aB-bB_2)Y + (cA_{21}-bA_2) + cB_{21}Y_{21} = aAX + cB_{21}Y_{21}
\end{align}
and this means that
\begin{align}
	aB = bB_2\com cA_{21} = b A_2\,.
\end{align}
We do not want to have any denominators, so when we solve for $a$, $b$ and $c$, we have a small solution
\begin{align}
	a = A_{21}B_2\com b= A_{21}B\com c = A_2B\,,
\end{align}
and these may or may not have common factors. In other words the elimination procedure yields
\begin{align}
\begin{split}
	&AA_{21}B_2 Q_1^{0} + A_2BB_{21}g_{21}Q_2^{0} =A_{21}B_2\beta-A_{21}B \beta_2 + A_2 B \beta_{21}\\
	&=(A_{21}B_2-A_{21}B+A_2B)\beta \,.
\end{split}	
\end{align}
Here we may think of a weighted sum operator too. Let
\begin{align}
	D(F) = A_{21}B_2 F - A_{21}B g_2(F) + A_2B g_{21}(F) \,.
\end{align}
Then we can rewrite the equation as
\begin{align}
	D(AX+BY) = AA_{21}B_2 Q_1^{0} + A_2BB_{21}g_{21}Q_2^{0} = \beta D(1)\,.
\end{align}
We may compute $D(1)$ for each case, but the expressions are somewhat complicated. $D(1)$ for each of $S11$, $S12$ and $S13$ will contain a square root like in $\sigma_1$. So do the terms $AA_{21}B_2$ and $A_2BB_{21}$. It is worth mentioning that we can convert the equation to one with no square roots by collecting terms involving the square root on one side and squaring. Then we have a nonlinear equation in $X$ and $Y$ with rational function coefficients. It is not clear how we could go further.
%=======================================
\section{Group of Order Eight}
%=======================================
Since the group $G$ has order 8, we only have two nonisomorphic models, but we will deal with the one with four group actions that preserve the admissibility of $(x_1,x_2)$. The actions are $1$, $ g_2$, $ g_{21}$, $ g_{212}$. The step polynomial is
\begin{align}
	\modelgraph 10011001 \,:\,S &= \ol x_1+x_1+\ol x_1x_2+x_1\ol x_2
\end{align}
This is the case covered in \cite{tabbara2014exact}. We will go through the same problem but from this new perspective. From the previous chapter we have
\begin{table}[H]
\caption{Reductions of group actions $g \in G$ on $Q_1^0$ and $Q_2^0$\\that preserve $1$-admissibility.}
\label{tbl:tableG6}
\centering
\begin{tabular}{ccc}
\toprule
	Sequence $L$ & $Q_1^0$ & $Q_2^0$ \\ 
\midrule
	$\emptyset$ & $\emptyset$ & $\emptyset$ \\
	2 &  2& $\emptyset$ \\
	21 &  2 & 21   \\
	212 &  212 & 21  \\
\bottomrule
\end{tabular}
\end{table}
We will (?)
Let $\beta = K_{12}\Qf{0,0} - \alpha = \alpha(b_1b_2\Qf{0,0}-1)$, $X = Q_1^0$ and $Y = Q_2^0$. Note that $\beta$ is left unchanged by group actions. We have the four equations
\begin{align}
	AX+BY&= \beta\,, \\
	A_2 X_2 + B_2 Y &=\beta\,,\\
	A_{21}X_2+ B_{21}Y_{21} &= \beta\,, \\
	A_{212} X_{212} + B_{212} Y_{21}&= \beta\,,
\end{align}
where we wish to eliminate variables $X$, $Y$ that contain $\sigma_1$ in their input. Those are $X_2$, $Y$ and $Y_{21}$, and we can eliminate them by first eliminating $Y$ and $Y_{21}$ from the first two and second two equations, and then eliminating $X_2$ from the resulting two equations. That is, we compute $B_2$ times the first equation minus $B$ times the second equation, and $B_{212}$ times the third equation minus $B_{21}$ from the last equation. In other words this is the same as if we apply the operators
\begin{align}
	D_2[B_2] &= B_2 - B g_2\,,\\
	- g_{212}D_2[B_2] &= g_{212}\,B g_2 - g_{212}\,B_2 =  B_{212} \,g_{21} - B_{21}\,  g_{212} 
\end{align}
to both sides of the master functional equation and substitute $x_1 = \sigma_1$. We obtain the equations
\begin{align}
	 AB_2 X - A_2 B X_2 &= \beta(B_2 - B)\,,\\
	  A_{21}B_{212} X_2 - A_{212} B_{21} X_{212} &= \beta(B_{212}-B_{21})\,.
\end{align}
To eliminate the $X_2$ terms we take $A_{21}B_{212}$ times the first equation and add $A_2 B$ times the second equation. We have
\begin{align}
	A_{21}B_{212}AB_2 X - A_2 BA_{212} B_{21} X_{212} = \beta(A_{21}B_{212}(B-B_2)+A_2 B(B_{212}-B_{21}))
\end{align}
We could recover the same result by applying the following operator $D$ to both sides of the master functional equation.
\begin{align}
\begin{split}
	D &=A_{21}B_{212} D_2[B_2] - A_2 B g_{212}D_2[B_2]\\
	   &= (A_{21}B_{212} - A_2 B g_{212})D_2[B_2]\\
	   &= (A_{21}B_{212} - g_{212}A_{21}B_{212})D_2[B_2]\\
	   &= D_{212}[A_{21}B_{212}]D_2[B_2]\,.
\end{split}
\end{align}
Starting from the master functional equation, we have The right-hand side
\begin{align}
	D(\beta) = \beta D_{212}[A_{21}B_{212}](B_2-B)
\end{align}
and the left-hand side
\begin{align}
\begin{split}
		 D(AX+BY) &= D(AX) \\
		&= D_{212}[A_{21}B_{212}]D_2[B_2](AX)\\
		&=D_{212}[A_{21}B_{212}](AB_2X - A_2BX_2)\\
		&=D_{212}[A_{21}B_{212}](AB_2X)\\
		&= D_{212}[AA_{21}B_2B_{212}](X)
\end{split}
\end{align}
We note that $A_{212} = A$, so we can actually factor out $A$ from the operator. Finally, we have the following theorem.
\begin{theorem} For the first model with group of order 8, we have
\begin{align}
	D_{212}[A_{21}B_2B_{212}]\left(Q_1^0\right) = \frac{\beta}{A} D_{212}[A_{21}B_{212}](B_2-B)
\end{align}
where $A_L = g_L(K_1)$, $B_L = g_L(K_2)$ and $\beta = \alpha (b_1b_2\Qf{0,0}-1)$, and the equality is valid only once we set $x_1 = \sigma_1$ in this expression.
\end{theorem}
If we went through the $\sigma_2$ substitution instead, we would have obtained an analogous statement. Just switch the $A$'s and the $B$'s, and the 1's and the 2's.

The interesting fact (as found in \cite{tabbara2014exact}) is that the coefficients in front of the $Q$'s become Laurent polynomials once we take into consideration the $\sigma_k$-substitutions, even though the $\sigma_k$'s contain square roots. The coefficient in front of $\beta$ can be reduced to a rational function plus a rational function times a square root after properly conjugating. The calculations are exceedingly complicated and would be too long to be done by hand. The square root part does not depend on $a_2$ in the first case, and $a_1$ in the second case. This is what allows the authors to find relationships between $\Qf{0,0}$ and expressions where one or more of $a_1$ and $a_2$ equal to 1.

%=======================================
\section{Highly-symmetric Two-dimensional Models}
%=======================================
We will continue further with the OKM for highly-symmetric models. There are only a few such models. If $x_k$ is in $S$, then so is $\ol x_k$. If one of the diagonal terms is in $S$, then so are the others. To have a group of order 4 we need to have both an $x_1$ and $x_2$ variable in $S$, so we have several nonisomorphic models. One with cartesian steps, one with diagonal steps, one with diagonal step and $x_1$,$\ol x_1$ and finally, the model with all steps. The notation is the same as in previous chapters. We derived the following equations
\begin{align}
	x_1 = \sigma_1:\quad\Delta_2(K2^\ast Q_1^0) &=\frac{1}{K_1} \left(\Delta_2(K_2^\ast K_{12})\Qf{0,0} - \alpha(K_2^\ast-K_2)\right)\,,\\
	x_2 = \sigma_2:\quad\Delta_1(K1^\bullet Q_2^0) &=\frac{1}{K_2} \left(\Delta_1(K_1^\bullet K_{12})\Qf{0,0} - \alpha(K_1^\bullet-K_1)\right)\,.
\end{align}
It is enough to only consider the analysis for one of the equations for now. Only the third model requires looking at both equations. We already know from the previous chapter that $\mu_2 = S_2^{\um 1}/S_2^1 = 1$ and 
\begin{align}
	\alpha(K_2^\ast -K_2)  &= \alpha t(\ol x_2 - x_2)S_2^{\um 1}(\sigma_1)\,,\\
	\Delta_2(K_2^\ast K_{12}) &= c_2t(\ol x_2-x_2)(c_1S_2^{\um 1}(\sigma_1) - \ol \sigma_1\Sf{\um 1,\um 1})\,.
\end{align}
To compute $\Delta_2(K2^\ast Q_1^0)$ we need to compute $K_2(\sigma_1,x_2)$ and $K_2^\ast(\sigma_1,x_2)$, which are given by
\begin{align}
	K_2 = c_2 -t\ol x_2S_2^{\um 1}(\sigma_1) \com K_2^\ast = c_2 -t x_2S_2^{\um 1}(\sigma_1)\,.
\end{align}
Therefore, we focus on computing $S_2^{\um 1}(\sigma)$ first.
\begin{align}
\begin{split}
	S_2^{\um 1}(\sigma_1) &= \Sf{0,1}+ \Sf{1,1}(\ol \sigma_1+\sigma_1)\\
	&= \Sf{0,1} + \Sf{1,1}\frac{1-tS_1^0}{t S_1^1}\\
	&= \frac{1}{tS_1^1}\left(\Sf{1,1}+t\left(\Sf{0,1}S_1^1-\Sf{1,1}S_1^0\right)\right)\\
	&= \frac{\Sf{1,1}+t\Sf{0,1}\Sf{1,0}}{t\left(\Sf{1,0}+\Sf{1,1}(\ol x_2+x_2)\right)}\,.
\end{split}
\end{align}
Therefore, the function $S_2^{\um 1}(\sigma_1)$ is a rational function of $x_2$.
\begin{table}[H]
\caption{Function $S_2^{- 1}(\sigma_1)$ for highly-symmetric models.}
\label{tbl:Symmetric}
\centering
\begin{tabular}{ccccc}
\toprule
	Model & $\Sf{0,1}$ & $\Sf{1,0}$ & $\Sf{1,1}$ & $S_2^{\um 1}(\sigma_1)$\\
\midrule
	\cartmodel & 1 & 1 & 0 & 1 \\
	\diagmodel & 0 & 0 & 1 & $\ol t /(\ol x_2+x_2)$\\
	\modelgraph 11100111 & 1 & 0 & 1 & $\ol t /(\ol x_2+x_2)$\\
	\modelgraph 10111101 & 0 & 1 & 1 & $\ol t /(\ol x_2+1+x_2)$\\
	\fullmodel & 1 & 1 & 1 & $(1+\ol t )/(\ol x_2+1+x_2)$\\
\bottomrule
\end{tabular}
\end{table}
\begin{lemma} For highly-symmetric models with $e = \Sf{1,1}$, $f = \Sf{1,0}$ and $g = \Sf{0,1}$, we have
\begin{align}
	\frac{\alpha(K_2^\ast -K_2)}{K_1}  &= \frac{\alpha (\ol x_2 - x_2)(e+tfg)}{S_1^1K_1}\,,\\
	\frac{\Delta_2(K_2^\ast K_{12})}{K_1} &= \frac{c_2(\ol x_2-x_2)(c_1(e+tfg)- et \ol \sigma_1S_1^1)}{S_1^1(c_1-t\ol \sigma_1S_1^1)}\,.
\end{align}
\end{lemma}
%=======================================
\subsection{Cartesian Model -- N, S, E, W}
%=======================================
In this case $S_2^{\um 1}(\sigma_1) = 1$, so a lot of the calculations will be simpler, and this is a good place to start further work. Here $K_{12} = c_1c_2$, so it behaves like $\alpha$ and we have
\begin{align}
	\Delta_2(K2^\ast Q_1^0) + \frac{\alpha\beta(K_2^\ast-K_2)}{K_1}=0\,,
\end{align}
where
\begin{align}
	\beta =  1-b_1b_2 Q_{1,2}^{0,0}
\end{align}
and
\begin{align}
	\frac{\alpha\beta(K_2^\ast -K_2)}{K_1}  &= \frac{\alpha\beta t (\ol x_2 - x_2)}{K_1} = \alpha\beta f(c_1,x_2)\,,
\end{align}
since $K_1$ does not depend on $c_2$ or $x_1$. The equation finally becomes
\begin{align}
	(c_2-tx_2)Q_1^0(x_2) - (c_2-t\ol x_2)Q_1^0(\ol x_2) + \alpha\beta f = 0\,.
\end{align}
We perform the extraction $[x_2^1]$ to get
\begin{align}
	c_2Q_{1,2}^{0,1}-t Q_{1,2}^{0,0} + \alpha\beta f_2^1 = 0
\end{align}
and the extraction $[b_2^0]$ to get
\begin{align}
	- t \wc Q_{1,2}^{0,0} +  \ol a_1f_2^1\,.
\end{align}
Using the second equation within the first one and simplifying yields
\begin{align}
	b_2 Q_{1,2}^{0,1} &= t(a_2 Q_{1,2}^{0,0} - \beta \wc Q_{1,2}^{0,0})\,,\\
	b_1 Q_{1,2}^{1,0} &= t(a_1 Q_{1,2}^{0,0} - \beta \wh Q_{1,2}^{0,0})
\end{align}
Here the second equation is obtained via the same method but with $\bullet$ instead of $\ast$. Using the same method we can also write
\begin{align}
	b_2 \Qf{1,1} = t(a_2 \Qf{1,0}-\beta \wc Q_{1,2}^{1,0})\,.
\end{align}
Multiplying by $b_1$ and using the above relations leads to
\begin{align}
\begin{split}
	b_1b_2 \Qf{1,1} &= t(a_2b_1\Qf{1,0}-\beta b_1 \Qc{1,0})\\
			      &=t a_2t(a_1\Qf{0,0}-\beta \Qh{0,0})-t\beta t(a_1\Qc{0,0}-\beta\Qch{0,0})\\	
			      &=t^2(a_1a_2\Qf{0,0}-a_2\beta\Qh{0,0}-a_1\beta\Qc{0,0}+\beta^2\Qch{0,0})\,.
\end{split}
\end{align}
Let
\begin{align}
	G = \Qf{0,0}, \quad\wh G = G_1, \quad\wc G = G_2\,.
\end{align}
From before we have
\begin{align}
	\Qf{0,0} = 1 + t a_1a_2(\Qf{1,0}+\Qf{0,1})\,.
\end{align}
Multiplying by $b_1b_2$ and using the above relations leads to
\begin{align}
\begin{split}
b_1b_2 G &= b_1b_2 + ta_1a_2(b_2b_1\Qf{1,0}+b_1b_2\Qf{0,1})\\
                 &=b_1b_2 + ta_1a_2(b_2t(a_1G-\beta G_1)+b_1t(a_2 G- \beta G_2)\\
                 &=b_1b_2 + t^2a_1a_2(b_2a_1G+b_1a_2G-b_2\beta G_1-b_1\beta G_2)\,.
\end{split}
\end{align}
Using the relation for $\beta$ and simplifying leads to
\begin{align}
	G = \frac{b_1b-2 - t^2 a_1a_2(b_2G_1+b_1G_2)}{b_1b_2 - t^2a_1a_2 (a_2b_1+a_1b_2+b_1b_2(b_2G_1+b_1G_2))}\,.
\end{align}
If we let $F = (b_2G_1+b_1G_2)$ and perform the long division with respect to $F$, we find that
\begin{align}
	\Qf{0,0} = \frac{1}{b_1b_2}\left(1-\frac{b_1b_2(1-b_1b_2)-t^2a_1a_2(a_1b_2+a_2b_1)}{b_1b_2+t^2a_1a_2(a_1b_2+a_2b_1+b_1b_2 F)}\right)\,.
\end{align}
Using arguments similar to those in the later section on diagonal steps we can find out more about the various $Q$'s. We have
\begin{align}
	b_2\Qf{m,n+1} &= t(a_2\Qf{m,n}-\beta \Qc{m,n})\,,\\
	b_1\Qf{m+1,n} &= t(a_1\Qf{m,n}-\beta \Qh{m,n})\,.
\end{align}
Summing the second set of equations multiplied by $x_2^n$, we find the equations, indexed by $m$,
\begin{align}
	b_1 Q_1^{m+1} &= t(a_1Q_1^m-\beta \wh Q_1^m)\,.
\end{align}
Now we sum these equations with respect to $x_1$ to find
\begin{align}
	b_1 (Q-Q_1^0) &= x_1t(a_1Q - \beta \wh Q)\,,
\end{align}
 and we can solve this equation for $Q_1^0$ as follows
\begin{align}
	b_1Q_1^0 &= (b_1-a_1x_1t)Q + x_1t\beta \wh Q\,.
\end{align}
We would have found an analogous expression for $b_2Q_2^0$ if we performed the procedure on the first set of equations by summing with respect to $x_1$ first and then with respect to $x_2$. We substitute these two expressions into the original functional equation, and after some algebra we find that
\begin{align}
	Q = \frac{(1-b_1b_2 G)(\alpha + tx_1\ol b_1K_1\wh Q+tx_2\ol b_2K_2 \wc Q)}{K-(1-\ol c_1x_1t)K_1-(1-\ol c_2x_2t)K_2}
\end{align}
We recognize some of the quantities as modified $K_1$ and $K_2$'s, so we can finally write down the full solution for $Q$.
\begin{theorem}The Cartesian model in two dimensions has generating functions satisfying the equations
\begin{align}
	\Qf{0,0} = G = \frac{1}{b_1b_2}\left(1-\frac{b_1b_2(1-b_1b_2)-t^2a_1a_2(a_1b_2+a_2b_1)}{b_1b_2+t^2a_1a_2(a_1b_2+a_2b_1+b_1b_2 F)}\right)\,,
\end{align}
where $F = (b_2\wh G+b_1\wc G)$,
\begin{align}
	b_1Q_1^0 &= (b_1-a_1x_1t)Q + x_1t\beta \wh Q\,,\\
	b_2Q_2^0 &= (b_2-a_2x_2t)Q + x_2t\beta \wc Q\,,
\end{align}
and 
\begin{align}
	Q = \frac{(1-b_1b_2 \Qf{0,0})(\alpha - \ol b_1 K_1 \wh K_1^\bullet \wh Q-\ol b_2K_2\wc K_2^\ast \wc Q)}{K - \ol c_1 K_1 K_1^\bullet - \ol c_2 K_2 K_2^\ast}\,.
\end{align}
\end{theorem}

This rather complicated expression shows that the generating function$Q$ is built from copies of $Q$ that have some of their marking variables set to 1.
%=======================================
\subsection{Diagonal Cartesian Model -- NE, NW, SE, SW}
%=======================================
The next case in terms of difficulty is when we have diagonal steps. This model is very similar to the one covered in \cite{wong2014confining} where they have two interacting walks confined in a semi-infinite strip. The models are different because there we have an extra term in the functional equations. Perhaps the difficulty in the current case is due to $S_2^{\um 1}(\sigma_1)$ being more complicated, and this is due to the diagonal terms (?). Also, $\Sf{\um 1, \um 1} \neq 1$, so $K_{12}$ is not trivial. Here
\begin{align}
	\frac{\alpha(K_2^\ast -K_2)}{K_1}  &= \frac{\alpha (\ol x_2 - x_2)}{K_1(\ol x_2+x_2)}\,,\\
	\frac{\Delta_2(K_2^\ast K_{12})}{K_1} &= \frac{c_2(\ol x_2-x_2)(c_1- t \ol \sigma_1S_1^1)}{S_1^1(c_1-t\ol \sigma_1S_1^1)} = \frac{c_2(\ol x_2-x_2)}{\ol x_2+x_2}\,.
\end{align}
The functional equation becomes
\begin{align}
	K_2^\ast Q_1 - K_2 Q_1^\ast -\frac{c_2(\ol x_2-x_2)}{\ol x_2+x_2}Q_{1,2}+\alpha\frac{\ol x_2-x_2}{(\ol x_2+x_2)K_1}= 0\,.
\end{align}
Surprisingly enough, the only term depending on $\sigma_1$, that remains in this expression, is $K_1$, so this will be of no consequence to the rest of the analysis. We just need to take care of the remaining terms. We will do this by multiplying the equation by $\ol x_2+x_2$. We have
\begin{align}
	R^\ast \exx= (\ol x_2+x_2) K_2^\ast &= c_2(\ol x_2+x_2)-x_2 = c_2\ol x_2- \ol a_2x_2\,,\\
	R = (\ol x_2+x_2) K_2 &= c_2(\ol x_2+x_2)-\ol x_2 = c_2x_2- \ol a_2\ol x_2\,.
\end{align}
There are a couple of relations here
\begin{align}
	R^\ast - R &= v(K_2^\ast-K_2) = \ol x_2-x_2\\
	\ol x_2R^\ast - x_2 R &= \ol x_2c_2v-1- x_2 c_2v+1 \\
				      & = c_2v(\ol x_2-x_2)
\end{align}
The equation becomes
\begin{align}
	(c_2\ol x_2- \ol a_2x_2)Q_1^0(x_2)+ ( \ol a_2\ol x_2-c_2x_2) Q_1^0(\ol x_2)+(\alpha \ol K_1 - c_2Q_{1,2}^{0,0})(\ol x_2-x_2)= 0
\end{align}
and it is ready for extraction. Let $\delta_0 =  \ol a_1(\ol x_2-x_2)\ol K_1$. We rewrite the equation as
\begin{align}\label{zero-order}
	R^\ast Q_1^0 -R Q_1^{0,\ast} - c_2(\ol x_2-x_2)\Qf{0,0} + \ol a_2\delta_0 = 0\,.
\end{align}
We will use the master equation
\begin{align}
	(1-tS)Q - (c_1-t\ol x_1S_1^{\um 1}) Q_1^0- (c_2-t\ol x_2S_2^{\um 1}) Q_2^0+ (c_1c_2-t\ol x_1\ol x_2S_{1,2}^{\um 1,\um 1})Q_{1,2}^{0,0}  = \ol a_1\ol a_2
\end{align}
to further our work here. For diagonal steps it is
\begin{align}
	(1-tS)Q - (c_1-tv\ol x_1) Q_1^0- (c_2-t\ol x_2S_2^{\um 1}) Q_2^0+ (c_1c_2-t\ol x_1\ol x_2)Q_{1,2}^{0,0}  = \ol a_1\ol a_2
\end{align}
extracting the $[x_1^0]$ coefficient yields
\begin{align}
	Q_1^0 =  \ol a_2+c_2\Qf{0,0} + a_1tvQ_1^1\,.
\end{align}
We can extract higher coefficients too. Let $k\geq 1$. Then the $[x_2^k]$ coefficient is
\begin{align}\begin{split}
	0 &= Q_1^k-tS_{1}^{\um 1}Q_1^{k+1}-tS_1^{0}Q_1^k-tS_1^1Q_1^{k-1}-c_2\Qf{k,0} \\ 
	   &\quad+ t\ol x_2(\Sf{\um 1,\um 1}\Qf{k+1,0}+\Sf{0,\um 1}\Qf{k,0}+\Sf{1,\um 1}\Qf{k-1,0})\\
	   &=Q_1^k-tvQ_1^{k+1}-tvQ_1^{k-1}-c_2\Qf{k,0} + t\ol x_2(\Qf{k+1,0}+\Qf{k-1,0})\,.
\end{split}\end{align}
We can rewrite the last one in a form more convenient for the next few steps
\begin{align}\label{formula}
	Q_1^k = tv(Q_1^{k+1}+Q_1^{k-1}) + c_2 \Qf{k,0} - t\ol x_2(\Qf{k+1,0}+\Qf{k-1,0})\,.
\end{align}
We will use induction to prove a nice formula, but for that, we first need to establish another base case because of this splitting of cases for the coefficient extraction. We have
\begin{align}
	Q_1^0 &=  \ol a_2+c_2\Qf{0,0} + a_1tvQ_1^1\,,\\
	Q_1^{0,\ast} \exx&=  \ol a_2+c_2\Qf{0,0} + a_1tvQ_1^{1,\ast}\,.
\end{align}
Combining them with equation (\ref{zero-order}) gives
\begin{align}
\begin{split}
        0 &= R^\ast Q_1^0 -R Q_1^{0,\ast} - c_2(\ol x_2-x_2)\Qf{0,0} + \ol a_2\delta_0 \\
	&= R^\ast (  \ol a_2+c_2\Qf{0,0} + a_1tvQ_1^1) -R ( \ol a_2+c_2\Qf{0,0} + a_1tvQ_1^{1,\ast}) \\
	&\quad- c_2(\ol x_2-x_2)\Qf{0,0} + \ol a_2\delta_0\\
	&=a_1tv(R^\ast Q_1^1-R Q_1^{1,\ast}) + (R^\ast-R)( \ol a_2+c_2\Qf{0,0}) \\
	&\quad- c_2(\ol x_2-x_2)\Qf{0,0} + \ol a_2\delta_0\\
	& =a_1tv(R^\ast Q_1^1-R Q_1^{1,\ast})+ \ol a_2(\ol x_2-x_2)+ \ol a_2\delta_0\\
	& =a_1tv(R^\ast Q_1^1-R Q_1^{1,\ast})+ \ol a_2((\ol x_2-x_2)+\delta_0)\,.
\end{split}
\end{align}
So we can write the equation
\begin{align}
	R^\ast Q_1^1-R Q_1^{1,\ast}+ \ol a_2\delta_1 = 0
\end{align}
where
\begin{align}
	\delta_1 = \frac{(\ol x_2-x_2)+\delta_0}{a_1tv}\,.
\end{align}
This equation together with equation (\ref{zero-order}) form the base cases for the following lemma
\begin{lemma} Let $m$ be an integer greater than zero. Then
\begin{align}\label{coeff-eqn}
	R^\ast Q_1^m -R Q_1^{m,\ast} - c_2(\ol x_2-x_2)\Qf{m,0} + \ol a_2\delta_m = 0
\end{align}
where $\delta_m$ is not a function of $a_2$.
\end{lemma}
\begin{proof}We show this by induction. We have the base cases $m = 0$ and $m=1$. The latter one fits the pattern because $\Qf{1,0} = 0$ since we cannot have walks ending on $(1,0)$. Now we assume that the formula is true for all cases up to and including $m$. Then we use formula (\ref{formula})
\begin{align}
\begin{split}
	& R^\ast Q_1^m -R Q_1^{m,\ast}\\
	=&tv(R^\ast Q_1^{m+1}-R Q_1^{m+1,\ast}) + tv(R^\ast Q_1^{m-1}-R Q_1^{m-1,\ast})\\
	   & + c_2(R^\ast-R)\Qf{m,0}-t(\ol x_2R^\ast-x_2 R)(\Qf{m+1,0}+\Qf{m-1,0})\\
	=& tv(R^\ast Q_1^{m+1}-R Q_1^{m+1,\ast}) + tv(c_2(\ol x_2-x_2)\Qf{m-1,0}- \ol a_2\delta_{m-1})\\
	   & + c_2(\ol x_2-x_2)\Qf{m,0}-c_2tv(\ol x_2-x_2)(\Qf{m+1,0}+\Qf{m-1,0})\\
	=& tv(R^\ast Q_1^{m+1}-R Q_1^{m+1,\ast})+ c_2(\ol x_2-x_2)\Qf{m,0}\\
	   &-c_2tv(\ol x_2-x_2)\Qf{m+1,0}) -  \ol a_2tv\delta_{m-1})\,.
\end{split}
\end{align}
Combining this with equation (\ref{coeff-eqn}) leads to
\begin{align}
	0 &= tv(R^\ast Q_1^{m+1}-R Q_1^{m+1,\ast})-c_2tv(\ol x_2-x_2)\Qf{m+1,0} -  \ol a_2tv\delta_{m-1} + \ol a_2\delta_m\,, \\
	0 & = R^\ast Q_1^{m+1}-R Q_1^{m+1,\ast}-c_2(\ol x_2-x_2)\Qf{m+1,0}+  \ol a_2\delta_{m+1}
\end{align}
where
\begin{align}
	\delta_{m+1} =\frac{\delta_m}{tv}-\delta_{m-1}\,.
\end{align}
This function $\delta_{m+1}$ is also independent from $a_2$ and this completes the proof!
\end{proof}
There is an equivalent statement for $Q_2^n$. We can expand $R^\ast$ and $R$ to a form we can use
\begin{align}
\begin{split}
	0 &=(c_2\ol x_2- \ol a_2x_2)Q_1^m(x_2) -( c_2x_2- \ol a_2\ol x_2)Q_1^m(\ol x_2) \\
	&\quad- c_2(\ol x_2-x_2)\Qf{m,0} + \ol a_2\delta_m \,.
\end{split}
\end{align}
We start by extracting the $[x_2^1]$ coefficient
\begin{align}
\begin{split}
	& c_2\Qf{m,2}- \ol a_2\Qf{m,0}-c_2\Qf{m,0} +c_2\Qf{m,0}+ [x_2^1] \ol a_2\delta_m = 0\\
	& c_2\Qf{m,2}- \ol a_2\Qf{m,0}+ [x_2^1] \ol a_2\delta_m = 0\,.
\end{split}
\end{align}
We can also extract the $[x_2^{n+1}]$ coefficient for $n\geq 2$
\begin{align}
 	c_2 Q_{1,2}^{m,n+2}- \ol a_2Q_{1,2}^{m,n}+[x_2^{n+1}] \ol a_2\delta_m= 0\,.
\end{align}
This formula also contains the $n = 0$ case, so we treat them together. We can set $a_2 = 1$ to get
\begin{align}
	-\wc Q_{1,2}^{m,n}+[x_2^{n+1}]\delta_m = 0\,.
\end{align}
Substituting this back into the above equation gives
\begin{align}
	b_2 Q_{1,2}^{m,n+2} = Q_{1,2}^{m,n} - \wc Q_{1,2}^{m,n}\,.
\end{align}
An equivalent formula can be derived if the process is followed through the use of $x_1$ instead of $x_2$ in the analysis.
\begin{corollary} For integers $m$ and $n$ greater than zero, we have the following relations
\begin{align}
	b_1 \Qf{m+2,n} &= Q_{1,2}^{m,n}-\wh Q_{1,2}^{m,n}\,,\\
	b_2 \Qf{m,n+2} &= Q_{1,2}^{m,n}-\wc Q_{1,2}^{m,n}\,.
\end{align}
\end{corollary}

We can sum up those equations to retrieve the following functional equations.
\begin{corollary} For integers $m$ greater than zero, we have the following relations
\begin{align}
	Q_1^m - \wh Q_1^m &= b_1 Q_1^{m+2}\,,\\
	\wc Q_1^m - \wch Q_1^m &= b_1 \wc Q_1^{m+2}\,,\\
	x_1^2(Q_2^n-\wh Q_2^n) &= b_1 (Q_2^n - Q_{1,2}^{1,n}x_1 - Q_{1,2}^{0,n})\,,\\
	x_1^2(\wc Q_2^n-\wch Q_2^n) &= b_1 (\wc Q_2^n - \wc Q_{1,2}^{1,n}x_1 - \wc Q_{1,2}^{0,n})\,,\\
	x_1^2(Q-\wh Q) &= b_1(Q-x_1Q_1^1-Q_1^0)\,,\\
	x_1^2(\wc Q-\wch Q) &= b_1(\wc Q-x_1\wc Q_1^1-\wc Q_1^0)\,.
\end{align}
\end{corollary}

Every other equation results from the previous one by either setting $a_1 = 1$ or $a_2 = 1$. We have an equivalent set of equations for the second coordinate --- just replace the subscript 1 with 2 and $m$ with $n$, to have consistent notation.

Now we can move onto the main result about the diagonal walks, that is, the function for labelled walks ending at $(0,0)$ is connected to the functions for semi-labelled and unlabelled walks ending on $(0,0)$. We can use the master equation to extract the $[x_1^0x_2^0]$ and $[x_1^1x_2^1]$ coefficients, which will produce the following two relations
\begin{align}
	\Qf{0,0} &= 1 + ta_1a_2 \Qf{1,1}\\
	\Qf{1,1} &= t(\Qf{0,0}+\Qf{2,0}+\Qf{0,2}+\Qf{2,2})\,.
\end{align}
We also have
\begin{align}
\begin{split}
	b_1b_2\Qf{2,2} &= b_1(\Qf{2,0}-\wc Q_{1,2}^{2,0})\\
	&= b_1 \Qf{2,0} - b_1\wc Q_{1,2}^{2,0}\\
	&= \Qf{0,0}-\wh Q_{1,2}^{0,0}-\wc Q_{1,2}^{0,0}+\wch Q_{1,2}^{0,0}\,,
\end{split}
\end{align}
together with
\begin{align}
	b_1b_2\Qf{2,0} &= b_2 \Qf{0,0} - b_2 \wh Q_{1,2}^{0,0}\,,\\
	b_1b_2\Qf{0,2} &= b_1 \Qf{0,0} - b_1 \wc Q_{1,2}^{0,0}\,.
\end{align}
We also have
\begin{align}
\begin{split}
b_1b_2\Qf{1,1} &= t(b_1b_2\Qf{0,0}+b_1b_2\Qf{2,0}+b_1b_2\Qf{0,2}+b_1b_2\Qf{2,2})\\
	&= t(b_1b_2\Qf{0,0}+ b_2 \Qf{0,0} - b_2 \wh Q_{1,2}^{0,0}+ b_1 \Qf{0,0} \\
	&\quad- b_1 \wc Q_{1,2}^{0,0}+\Qf{0,0}-\wh Q_{1,2}^{0,0}-\wc Q_{1,2}^{0,0}+\wch Q_{1,2}^{0,0})\\
	&= t((b_1b_2+b_1+b_2+1)\Qf{0,0}-(b_2+1)\wh Q_{1,2}^{0,0}\\
	&\quad-(b_1+1)\wc Q_{1,2}^{0,0}+\wch Q_{1,2}^{0,0})\\
	&=t(a_1a_2\Qf{0,0}-a_2\wh Q_{1,2}^{0,0}-a_1\wc Q_{1,2}^{0,0}+\wch Q_{1,2}^{0,0})\,,
\end{split}
\end{align}
We can finally write
\begin{align}
\begin{split}
	b_1b_2\Qf{0,0} &= b_1b_2 + ta_1a_2b_1b_2\Qf{1,1}\\
		& =  b_1b_2 + t^2a_1a_2(a_1a_2\Qf{0,0}-a_2\wh Q_{1,2}^{0,0}-a_1\wc Q_{1,2}^{0,0}+\wch Q_{1,2}^{0,0})\,.
\end{split}
\end{align}
Let $G(b_1,b_2)$ stand for $Q(0,0,b_1,b_2)$ and $G_k$ has $b_k = 0$. Then we can write it all as
\begin{align}
	b_1b_2 G = b_1b_2 + t^2a_1a_2(a_1a_2G-a_2G_1-a_1G_2+G_{1,2})\,,
\end{align}
or written as a linear combination,
\begin{align}
	(b_1b_2-t^2a_1^2a_2^2)G+t^2a_1a_2^2G_1 + t^2a_1^2a_2G_2 - t^2a_1a_2G_{1,2} = b_1b_2\,.
\end{align}
We can divide by $a_1a_2$ and obtain the equation
\begin{align}
	(c_1c_2-a_1a_2t^2)G + a_2t^2G_1+a_1t^2G_2-t^2G_{1,2}=c_1c_2\,,
\end{align}
To conclude, we have the following result.

\begin{theorem} Let $G = \Qf{0,0}$ be the generating function for walks ending at the origin in the Diagonal Cartesian model. Let $f(a_1,a_2,t) = c_1c_2-a_1a_2t^2$ and let $B = B_1B_2$. Then $B(fG) = c_1c_2$.
\end{theorem}
We can also solve for $Q_k^j$ and other terms, in terms of various $[b_k^l]$ extractions, so overall, it should be possible to find a relationship of the form $B(FQ) = 0$.

%=======================================
\subsection{The Full Model}
%=======================================
We will finally take a look at the model with all the steps.Here we have the most complicated functions. Here
\begin{align}
	\frac{\alpha(K_2^\ast -K_2)}{K_1}  &= \frac{\alpha (\ol x_2 - x_2)(1+t)}{S_1^1K_1}\,,\\
	\frac{\Delta_2(K_2^\ast K_{12})}{K_1} &= \frac{c_2(\ol x_2-x_2)(c_1(1+t)- t \ol \sigma_1S_1^1)}{S_1^1(c_1-t\ol \sigma_1S_1^1)}\\
	&=\frac{c_2(\ol x_2-x_2))}{S_1^1}\left(1-\frac{t}{c_1-t\ol \sigma_1S_1^1}\right)\,.
\end{align}
We will need to do some more work to find a representation for that last fraction. For that we need the $\sigma_1$ root, which is
\begin{align}
\begin{split}
	\sigma_1 &= \frac{(1-tS_1^0)-\sqrt{(1-tS_1^0)^2-4t^2(S_1^1)^2}}{2tS_1^1} \\
		&= \frac{u-\sqrt{u^2-4w^2}}{2w}\\
\end{split}\\
	&\ol \sigma_1 = \frac{u+\sqrt{u^2-4w^2}}{2w}\,.
\end{align}
Then we have
\begin{align}
\begin{split}
	\frac{t}{c_1-\ol \sigma_1w} &= \frac{2t}{2c_1-u-\sqrt{u^2-4w^2})} \\
	&= \frac{2t(2c_1-u+\sqrt{u^2-4w^2})}{(2c_1-u)^2-(u^2-4w^2)}\\
	&=  \frac{t(2c_1-u+\sqrt{u^2-4w^2})}{2(c_1^2-c_1u+w^2)}
\end{split}
\end{align}
Let $v = \ol x_2+x_2$. Then the denominator has
\begin{align}
\begin{split}
	c_1^2-c_1u+w^2 &= c_1^2-c_1(1-tS_1^0)+t^2(S_1^1)^2\\&=\ol a_1^2b_1 - \ol a_1b_1vt+(1+v)^2 t^2
\end{split}
\end{align}
and the numerator has
\begin{align}
\begin{split}
	u^2-4w^2 &= (1-tv)^2-4t^2(1+v)^2 \\&= 1-2tv-(3v+2)(v+2)t^2 \\&= (1-(3v+2)t)(1+(v+2)t)\,.
\end{split}
\end{align}
We can find the first few terms of the square root.
\begin{align}
\begin{split}
	\sqrt{u^2-4w^2} = 1 - vt-2(1+v)^2 t^2-2v(1+v)^2t^3+\cdots
\end{split}
\end{align}
We can give a series representation of the square root via the series
\begin{align}
	\sqrt{1-z} = \sum_{n=0}^\infty \frac{(2n)!}{(1-2n)(n!)^2 4^n}z^n\,.
\end{align}
That is,
\begin{align}
\begin{split}
	&\sqrt{(1+(v+2)t)(1-(3v+2)t)} \\=& \left(\sum_{n=0}^\infty \frac{(-1)^n(2n)!(v+2)^n t^n}{(1-2n)(n!)^2 4^n}\right)\left(\sum_{n=0}^\infty \frac{(2n)!(3v+2)^nt^n}{(1-2n)(n!)^2 4^n}\right)\\
	=& \sum_{n=0}^\infty \frac{t^n}{4^n}\sum_{k=0}^n \frac{(-1)^k(2k)!(2n-2k)!}{(1-2k)(1-2n+2k)k!^2 (n-k)!^2}(v+2)^k(3v+2)^{n-k}\,,
\end{split}
\end{align}
However, we are not sure if a closed form exists for these coefficients. If we want to extract $x_2$ coefficients from this expression, we would not have an easy time (?). Now, $K_2$ and $K_2^\ast$ have $1+v$ in their denominators, so we can multiply the whole functional equation by $1+v$ to clear out these denominator. We have
\begin{align}
	(1+v)K_2 &= c_2(1+v)-\ol x_2(1+t) = (c_2-1-t)\ol x_2+c_2+c_2x_2\,,\\
	(1+v)K_2^\ast &= c_2(1+v)-x_2(1+t) = c_2\ol x_2+c_2+(c_2-1-t)x_2\,,
\end{align}
and
\begin{align}
	(1+v)\frac{\alpha(K_2^\ast -K_2)}{K_1} = &\frac{\alpha (\ol x_2 - x_2)(1+t)}{K_1} = \alpha g(a_1,x_2,t)\,,
\end{align}
\begin{align}
\begin{split}
	(1+v)\frac{\Delta_2(K_2^\ast K_{12})}{K_1}=&c_2(\ol x_2-x_2)\left(1-t\frac{c_1-1+tv+\sqrt{(1+(v+2)t)(1-(3v+2)t)}}{c_1^2-c_1+c_1tv+(1+v)^2t^2}\right)\,,
\end{split}
\end{align}
and let the latter be $c_2(\ol x_2-x_2)f(a_1,x_2,t)$. Then if we let $A = c_2-(1+t)$, we can write the funtional equation as
\begin{align}
	(c_2\ol x_2+c_2+Ax_2)Q_1^0(x_2) - (A\ol x_2+c_2+c_2 x_2)Q_1^0(\ol x_2) = c_2(\ol x_2-x_2)f-\alpha g\,.
\end{align}
As before, we extract the $[x_2]$ coefficient from both sides to find that
\begin{align}
	 c_2\Qf{0,2}+c_2\Qf{0,1}-(1+t)\Qf{0,0} = -\alpha g_2^1 + c_2(f_2^2-f_2^0)\Qf{0,0}\,,
\end{align}
and then set $a_2 = 1$ on both sides to find that
\begin{align}
	-(1+t)\Qc{0,0} = -\ol a_1 g_2^1\,.
\end{align}
When we finally substitute this quantity back into the equation, we can simplify it to
\begin{align}
	b_2(\Qf{0,2}+\Qf{0,1}) = (b_2(f_2^2-f_2^0)-a_2(1+t))\Qf{0,0}-(1+t)\Qc{0,0}\,.
\end{align}
The question is can we find the quantity $f_2^2-f_2^0$. It appears that
\begin{align}
	-(f_2^2-f_2^0) =1+a_1t+\frac{a_1^2}{b_1}(a_1^2-a_1-4)t^3 +\frac{4a_1^2}{b_1^2}(a_1^3-3a_1^2+1)t^4+\cdots\,.
\end{align}
Unfortunately we have not been able to go further with this case. One could try making some assumptions for the form of equations like these for $\Qf{m,n}$, $\Qc{m,n}$ and $\Qh{m,n}$, but even then it is not clear if we can find a relation between $\Qf{0,0}$, $\Qc{0,0}$, $\Qh{0,0}$ and $\Qch{0,0}$. We would have 6 equations like the one above plus 6 more where $b_i = 0$ for the $b_i$ not present. Then one may use the 4 regular coefficient extractions from the master functional equation for $\Qf{m,n}$ and $m,n\leq 1$ plus 12 more where we set one or more of $b_1,b_2$ to zero. This brings us to a grand total of 28 equations in 36 unknown variables, of which we need to eliminate 32. It is not entirely clear if we may do that.

\endinput